% [原文 p.125]
\par\medskip\noindent\textbf{3.32}\quad 本节余下部分说明纲领的哪些部分可由已确立的结果推出，哪些部分仍待证明。假设 $f_0:X\to S$ 满足条件 3.21(1)--(2) 之一，下面逐点详细讨论。

\begin{enumerate}[label=(\arabic*)]
\item 我们用到了比定义 2.8 更一般的翻转概念。一般定义在定义 3.33 中给出。
\item 本节的设定中需要用到的一般形式的锥定理，在定理 3.35 中证明。
\item 命题 3.36 与命题 3.37 证明 Q-分解性条件在 MMP 过程中保持不变。
\item 推论 3.44 证明 dlt 条件在 MMP 过程中保持不变。更一般地，也处理了终端、标准、klt、lc 对（推论 3.42、推论 3.43）。
\item 研究极小模型纲领的结局，有助于确立对的一般的相对极小模型定义。最终会看出该定义相当微妙，这放在下一节。
\item 3.45 讨论 MMP 是否终止的问题。一般情形尚未解决。后续章节处理若干特殊情形。
\item 最大的未解决问题是翻转的存在性。本书余下部分本质上用于解决该问题的若干情形。\hfill$\square$
\end{enumerate}
% [Tutor Note: 第 (7) 条点出全书主线：翻转的存在性是三维 MMP 的唯一缺口。三维终端情形由森重文 [Mor88] 解决，本书第 6 章即为此而设。]

\begin{definition}{3.33}
设 $X$ 是正规概形（或正规解析流形），$D$ 是 $X$ 上的 Q-除子，且 $K_X+D$ 为 Q-Cartier。眼下不假设 $D$ 有效。

$(K+D)$-翻转收缩（$(K+D)$-flipping contraction）是指满足下述条件的、到正规射影概形（或正规解析流形）$Y$ 的固有双有理态射 $f:X\to Y$：$\operatorname{Ex}(f)$ 在 $X$ 内的余维数 $\ge 2$，且 $-(K_X+D)$ 为 $f$-丰富。

正规概形（或正规解析流形）$X^{+}$ 连同固有双有理态射 $f^{+}:X^{+}\to Y$，当下列条件满足时称为 $f$ 的（$K+D$）-翻转（$(K+D)$-flip）：
% [原文 p.126]
\begin{enumerate}[label=(\arabic*)]
\item $K_{X^{+}}+D^{+}$ 为 Q-Cartier，其中 $D^{+}$ 是 $D$ 在 $X^{+}$ 上的双有理变换；
\item $K_{X^{+}}+D^{+}$ 为 $f^{+}$-丰富；
\item $\operatorname{Ex}(f^{+})$ 在 $X^{+}$ 内的余维数 $\ge 2$。
\end{enumerate}
这里是记号的滥用：双有理映射 $\varphi=(f^{+})^{-1}\circ f:X\dashrightarrow X^{+}$ 也称为（$K+D$）-翻转。推论 6.4 将证明（$K+D$）-翻转若存在则唯一，因此主要问题是存在性。（$K+D$）-翻转给出如下交换图：

\begin{center}
\begin{tikzcd}
X \arrow[rr, dashed, "\varphi"] \arrow[dr, "f"'] & & X^{+} \arrow[dl, "f^{+}"] \\
& Y &
\end{tikzcd}
\end{center}

其中 $-(K_X+D)$ 为 $f$-丰富，$K_{X^{+}}+D^{+}$ 为 $f^{+}$-丰富。

文献中对术语的使用并不一致，上述操作有时也称为 $D$-翻转。

特别地，当 $\dim X=3$ 时，曲线 $C\subset\operatorname{Ex}(f)$ 称为前翻转曲线（flipping curve），曲线 $C^{+}\subset\operatorname{Ex}(f^{+})$ 称为后翻转曲线（flipped curve）。
\end{definition}
% [Tutor Note: 翻转收缩只要求 $\operatorname{Ex}(f)$ 余维数至少 2，即它是小收缩；且此处 $D$ 不必有效，故这一定义可直接沿 MMP 使用。]

MMP 过程中遇到的翻转还有其他有用的性质。下述条件往往十分便利。

\begin{definition}{3.34}
固有双有理态射 $f:X\to Y$ 在满足下述条件时称为极端的（extremal）或射线型的（ray type）。
\begin{enumerate}[label=(\arabic*)]
\item $X$ 是 Q-分解的。
\item 若 $B,B'$ 为 Cartier 除子，则存在不全为零的 $a,a'\in\Z$，使得 $aB\sim_f a'B'$。
\end{enumerate}
一般通行的是“极端的”这一称法，但按此说法，极锥面的收缩就不是极端的，诸如此类，颇为奇怪。本书主要使用“射线型”。

假设 $(X,\Delta)$ 为 klt，$X$ 为 Q-分解的。设 $f:X\to Y$ 是 $(K_X+\Delta)$-负的极射线 $R$ 的收缩，任意选取 $[C]\in R$。
% [原文 p.127]
选取不全为零的 $a,a'\in\Z$，使得 $a(B\cdot C)=a'(B'\cdot C)$。则由定理 3.25(4) 得 $aB\sim_f a'B'$。故 $f$ 是射线型的。
\end{definition}
% [Tutor Note: 射线型是相对 Picard 群秩 1 的条件，与“极射线”无关；改称“射线型”正是为了避免“极端的”与凸几何用语混淆。]

\begin{theoremp}{3.35}
相对锥定理 3.25 对 Q-分解的 dlt 对 $(X,\Delta)$ 与射影态射 $f:X\to S$ 成立。
\end{theoremp}

\begin{proof}
首先证明即使 $(X,\Delta)$ 为 dlt，定理 3.25(2) 也成立。选取 $0<\delta\ll1$，使得 $H':=\delta\Delta+\eps H$ 为 $f$-丰富。由命题 2.43，$(X,(1-\delta)\Delta)$ 为 klt，故由定理 3.25(2)，
\[
\NE(X/S) = \NE(X/S)_{K_X+(1-\delta)\Delta+H'\ge0} + \sum_{\text{有限}} \R_{\ge0}[C_j]
\]
成立，这等价于
\[
\NE(X/S) = \NE(X/S)_{K_X+\Delta+\eps H\ge0} + \sum_{\text{有限}} \R_{\ge0}[C_j].
\]
正如定理 3.15 第 5 阶段中注意到的，这在 dlt 情形意味着定理 3.25(1)。

套用这一论证，可以证明任意 $(K_X+\Delta)$-负的相对极射线，当 $0<\delta\ll1$ 时也是 $(K_X+(1-\delta)\Delta)$-负的相对极射线。$(X,(1-\delta)\Delta)$ 为 klt，故由定理 3.25(3)--(4) 的 klt 情形可推出 Q-分解 dlt 情形的相应断言。
\end{proof}
% [Tutor Note: 证明的关键是扰动：把 dlt 的 $\Delta$ 换成 $(1-\delta)\Delta+\eps H$，既保持负性又化为 klt，从而把已知的 klt 锥定理搬过来。]

利用推论 3.17 与推论 3.18，同样的论证证明 Q-分解性在 dlt 情形下也保持不变。

\begin{proposition}{3.36}
设 $(X,\Delta)$ 为射影的 Q-分解 dlt 对，$g_R:X\to Y$ 为 $(K_X+\Delta)$-负的极射线的收缩。假设 $g_R$ 是除子收缩或 Fano 收缩。则
\begin{enumerate}[label=(\arabic*)]
\item $Y$ 是 Q-分解的；
\item $\rho(Y)=\rho(X)-1$。
\end{enumerate}
\end{proposition}

Q-分解性在翻转下也保持不变。

\begin{propositionp}{3.37}
设 $(X,\Delta)$ 为射影的 Q-分解 dlt 对，$g_R:X\to Y$ 为 $(K_X+\Delta)$-负的极射线的翻转收缩，$g_R^{+}:X^{+}\to Y$ 为其翻转。则
% [原文 p.128]
\begin{enumerate}[label=(\arabic*)]
\item $X^{+}$ 是 Q-分解的；
\item $\rho(X^{+})=\rho(X)$。
\end{enumerate}
\end{propositionp}

\begin{proof}
因 $\varphi:X\dashrightarrow X^{+}$ 在余维数 1 处为同构，它给出 $X$ 的 Weil 除子群与 $X^{+}$ 的 Weil 除子群之间的自然同构。设 $D^{+}$ 是 $X^{+}$ 上的 Weil 除子，$D$ 是它在 $X$ 上的双有理变换。取适当的有理数 $r$，使得 $(R\cdot(D+r(K_X+\Delta)))=0$。选取 $m\in\Z_{>0}$，使 $m(D+r(K_X+\Delta))$ 为 Cartier 除子。由推论 3.17，存在 Cartier 除子 $D_Y$，满足 $m(D+r(K_X+\Delta))\sim g_R^*D_Y$。于是
\[
mD^{+} = m\varphi_*D \sim (g_R^{+})^*D_Y-(mr)(K_{X^{+}}+\Delta^{+})
\]
为 Q-Cartier。
\end{proof}

下述一般结果用于在几种不同的设定中比较差异系数。

\begin{lemmap}{3.38}
考虑下述交换图。
\begin{center}
\begin{tikzcd}
X \arrow[rr, dashed, "\varphi"] \arrow[dr, "f"'] & & X' \arrow[dl, "f'"] \\
& Y &
\end{tikzcd}
\end{center}
其中 $X,X',Y$ 为正规簇，$f,f'$ 为固有双有理态射。设 $\Delta$（相应地 $\Delta'$）是 $X$（相应地 $X'$）上的 Q-除子。假设以下 3 个条件：
\begin{enumerate}[label=(\arabic*)]
\item $f_*\Delta=f'_*\Delta'$；
\item $-(K_X+\Delta)$ 为 Q-Cartier 且为 $f$-nef；
\item $K_{X'}+\Delta'$ 为 Q-Cartier 且为 $f'$-nef。
\end{enumerate}
则对 $Y$ 上方任意的例外除子 $E$，有
\[
a(E,X,\Delta)\le a(E,X',\Delta').
\]
若下列条件之一满足，则上述不等式为严格不等号。
\begin{enumerate}[label=(\arabic*),start=4]
\item $-(K_X+\Delta)$ 为 $f$-丰富，且 $f$ 在 $\operatorname{center}_Y E$ 的泛点处不是同构；
\item $(K_{X'}+\Delta')$ 为 $f'$-丰富，且 $f'$ 在 $\operatorname{center}_Y E$ 的泛点处不是同构。
\end{enumerate}
\end{lemmap}
% [Tutor Note: 条件 (1) 保证 $H$ 是例外除子的线性组合，否则无法用引理 3.39 的符号结论；这是比较差异系数的标准套路。]

% [原文 p.129]
\begin{proof}
设 $Z$ 为正规簇，$g:Z\to X$ 与 $g':Z\to X'$ 为双有理态射，并假设 $\operatorname{center}_Z E$ 是除子。记 $h:=f\circ g=f'\circ g'$。取 $m>0$ 为充分大、且能被所需各整数整除的整数，则有如下线性等价：
\[
-m\Bigl(K_Z-\sum a(E_i,X,\Delta)E_i\Bigr)\sim -mg^*(K_X+\Delta),
\]
\[
m\Bigl(K_Z-\sum a(E_i,X',\Delta')E_i\Bigr)\sim m(g')^*(K_{X'}+\Delta').
\]
两式相加可知
\[
H:=\sum\bigl(a(E_i,X,\Delta)-a(E_i,X',\Delta')\bigr)E_i
\]
为 $h$-nef，且由假设 (1) 它是 $h$-例外除子的线性组合。由引理 3.39，其所有系数 $\le0$。进一步，若 $H$ 在 $\operatorname{center}_Y E$ 的泛点上不是数值 $h$-平凡的，则 $H$ 中 $E$ 的系数 $<0$。
\end{proof}

下述引理在许多情形下有用。

\begin{lemmap}{3.39}（nef 例外除子引理）
设 $h:Z\to Y$ 是正规簇之间的固有双有理态射。设 $-B$ 是 $Z$ 上的 $h$-nef Q-Cartier Q-除子。则以下成立。
\begin{enumerate}[label=(\arabic*)]
\item $B$ 有效当且仅当 $h_*B$ 有效。
\item 假设 $B$ 有效。则对每个 $y\in Y$，或者 $h^{-1}(y)\subset\operatorname{Supp}B$，或者 $h^{-1}(y)\cap\operatorname{Supp}B=\varnothing$。
\end{enumerate}
\end{lemmap}

\begin{proof}
若 $B$ 有效，则显然 $h_*B$ 有效。问题在于证明逆命题。

由 Chow 引理，存在固有双有理态射 $p:Z'\to Z$，使 $Z'\to Y$ 为射影的。注意 $B$ 有效当且仅当 $p^*B$ 有效。因此可以一开始就假设 $h$ 是射影的。也可以假设 $Y$ 是仿射的。把 $B$ 写成 $B=\sum B^k$，其中 $B^k$ 是 $B$ 的满足 $h(B_i)$ 在 $Y$ 内余维数为 $k$ 的所有不可约分支 $B_i$ 之和。

首先处理 $\dim Y=2$ 的情形。此时 $B=B^2+B^1$，且 $B^1$ 也是 $h$-nef。故 $-B^2$ 是例外曲线的 $h$-nef 线性组合。如前，$B^2$ 的有效性问题可归约到任意奇点消解 $U\to Z$ 的情形。由下面两个结果，$B^2$ 是有效的。第一个结果是 Hodge 指数定理的特殊情形。
% [Tutor Note: 原文此后转入引理 3.40 与 3.41，再回到本证明；阅读时应把这两条当作证明中插入的引理，而不是独立的平行结论。]

% [原文 p.130]
\begin{lemmap}{3.40}
设 $f:Y\to X$ 是正规曲面的奇点消解，$E_i$ 为例外曲线。则相交矩阵 $(E_i\cdot E_j)$ 是负定的。
\end{lemmap}

\begin{proof}
首先证明 $X$ 为射影的情形。设 $D=\sum e_iE_i$ 是例外曲线的非零线性组合，且假设 $(D^2)\ge0$。先考虑 $D$ 有效的情形。取 $Y$ 上的丰富除子 $H$，使得 $H-K_Y$ 为丰富。由 Serre 对偶，$H^2(Y,\mathcal{O}_Y(nD+H))=0$。又
\[
\bigl((nD+H)\cdot(nD+H-K_Y)\bigr)\ge\bigl((nD+H)\cdot(nD+H)\bigr)\ge n(D\cdot H).
\]
故如引理 1.20 所见，由曲面情形的 Riemann--Roch 定理，$n\to\infty$ 时 $h^0(Y,\mathcal{O}_Y(nD+H))\to\infty$。另一方面，
\[
H^0(Y,\mathcal{O}_Y(nD+H))\subset H^0(X,\mathcal{O}_X(f_*(nD+H)))=H^0(X,\mathcal{O}_X(f_*H)),
\]
矛盾。若 $D$ 不是有效的，考虑无公共分支的两个有效除子之差 $D=D_+-D_-$，则 $(D^2)\le(D_+^2)+(D_-^2)<0$。

$X$ 为拟射影的情形只需紧化。$X$ 为解析曲面的情形也成立，其证明委于 [Gra62, p.367]。
\end{proof}

\begin{lemmap}{3.41}
设 $U$ 为非奇异曲面，$C=\cup C_i$ 是 $U$ 上固有曲线的集合。假设相交矩阵 $(C_i\cdot C_j)$ 负定。设 $A=\sum a_iC_i$ 是曲线 $C_i$ 的 $\R$-线性组合。假设对每个 $j$ 都有 $(A\cdot C_j)\ge0$。则
\begin{enumerate}[label=(\arabic*)]
\item 对每个 $i$ 有 $a_i\le0$。
\item 若 $C$ 连通，则或者对每个 $i$ 有 $a_i=0$，或者对每个 $i$ 有 $a_i<0$。
\end{enumerate}
\end{lemmap}

\begin{proof}
把 $A$ 表示为无公共分支的两个非负线性组合之差 $A=A^{+}-A^{-}$。

假设 $A^{+}\ne0$。矩阵 $(C_i\cdot C_j)$ 负定，故 $(A^{+}\cdot A^{+})<0$。因此存在曲线 $C_i\subset\operatorname{Supp}A^{+}$，使 $(C_i\cdot A^{+})<0$。由于 $C_i\not\subset\operatorname{Supp}A^{-}$，有 $(C_i\cdot A^{-})\ge0$。于是 $(C_i\cdot A)<0$，矛盾。(1) 获证。

其次假设 $C$ 连通，且 $0\ne\operatorname{Supp}A^{-}$ 并且 $\operatorname{Supp}A^{-}\ne\operatorname{Supp}C$。此时存在不含于 $\operatorname{Supp}A^{-}$ 而与 $\operatorname{Supp}A^{-}$ 相交的曲线 $C_i$。
% [原文 p.131]
于是 $(C_i\cdot A)=-(C_i\cdot A^{-})<0$，矛盾。
\end{proof}

回到引理 3.39 的证明。设 $S\subset Y$ 是 $n-2$ 个一般超曲面的完全交，$T:=h^{-1}(S)$。此时 $h:T\to S$ 是正规曲面之间的双有理态射。$B|_T=B^2|_T+B^1|_T$，且由二维情形的结果，$B^2$ 有效。

设 $H\subset X$ 是充分丰富的一般除子，令 $A=B|_H$。此时 $-A$ 为 $h$-nef，$A^i=B^{i+1}|_H$（$i\ge2$），$A^1=B^1|_H+B^2|_H$。由假设 $B^1$ 有效，且已证 $B^2$ 有效，故 $A^1$ 也有效。对维数归纳，$A$ 有效，从而 $B$ 也有效。
% [Tutor Note: 原文作 $H\subset X$，疑为 $H\subset Z$ 之误；因 $B$ 定义在 $Z$ 上，只有 $H\subset Z$ 时 $A=B|_H$ 才有意义。此处照抄原文，特此说明。]

对 $y\in Y$，$h^{-1}(y)$ 连通（Zariski 主定理，[Har77, V.5.2]）。因此若 $h^{-1}(y)$ 与 $\operatorname{Supp}B$ 相交但不被其包含，则存在不可约曲线 $C\subset h^{-1}(y)$，使 $(C\cdot B)>0$。由于 $-B$ 为 $h$-nef，这不可能。
\end{proof}

引理 3.38 的第一个推论是：极小模型纲领不会使奇点变坏。

\begin{corollary*}{3.42}
设 $(X,\Delta)$ 为对，$\varphi:X\dashrightarrow X^{+}$ 为 $(K_X+\Delta)$-翻转。若 $(X,\Delta)$ 为终端（或标准、klt、lc），则 $(X^{+},\Delta^{+})$ 也为终端（或标准、klt、lc）。
\end{corollary*}

\begin{corollaryp}{3.43}
设 $(X,\Delta)$ 为对，$f:X\to Y$ 为 $(K_X+\Delta)$-负的极射线的除子收缩，$E$ 为例外除子。则以下成立。
\begin{enumerate}[label=(\arabic*)]
\item 若 $(X,\Delta)$ 为 klt（或 lc），则 $(Y,f_*\Delta)$ 同样如此。
\item 若 $(X,\Delta)$ 为终端（或标准）且 $E\not\subset\operatorname{Supp}\Delta$，则 $(Y,f_*\Delta)$ 也为终端（或标准）。
\item 若 $X$ 为终端（或标准）且 $\Delta=0$，则 $Y$ 也为终端（或标准）。
\end{enumerate}
\end{corollaryp}

\begin{proof}
对 $X'=X^{+}$ 应用引理 3.38，即得推论 3.42。对推论 3.43 则取 $X'=Y$。

终端性与标准性的假设只用例外除子的差异系数来定义。$E$ 在 $Y$ 上是例外的，但在 $X$ 上不是例外的。因此必须对 $E$ 单独处理，这就是 (2) 的表述的由来。
\end{proof}
% [Tutor Note: MMP 每一步都不使奇点类别变坏；但终端、标准只由例外除子定义，故除子收缩时必须单独检查被收缩的 $E$，这正是 (2) 中 $E\not\subset\operatorname{Supp}\Delta$ 的由来。]

% [原文 p.132]
\begin{corollaryp}{3.44}
设 $(X,\Delta)$ 为 dlt 对。设 $g:X\dashrightarrow X'$ 是 $(K_X+\Delta)$-负的极射线的除子收缩或 $(K_X+\Delta)$-翻转。令 $\Delta':=g_*\Delta$。则 $(X',\Delta')$ 也为 dlt。
\end{corollaryp}

\begin{proof}
按定义 2.37 取 $Z\subset X$，令 $Z':=g(Z)\cup\operatorname{Ex}(g^{-1})$。这里记号有滥用：$\operatorname{Ex}(g^{-1})$ 表示 $X'$ 中使 $g^{-1}$ 不是双正则的点全体。$X'\setminus Z'$ 与 $X\setminus Z$ 的开子集同构，故 $X'\setminus Z'$ 非奇异，且 $\Delta'|_{X'\setminus Z'}$ 是 snc 除子。

设 $E$ 是 $X'$ 上方满足 $\operatorname{center}_{X'}E\subset Z'$ 的例外除子，则 $\operatorname{center}_XE\subset Z\cup\operatorname{Ex}(g)$。故
\[
a(E,X',\Delta')\ge a(E,X,\Delta)\ge-1.
\]
上式中，若 $\operatorname{center}_XE\subset Z$，则由 dlt 的定义第二个不等式为严格不等号；若 $\operatorname{center}_XE\subset\operatorname{Ex}(g)$，则由引理 3.38 第一个不等式为严格不等号。
\end{proof}

\par\medskip\noindent\textbf{3.45}（MMP 的终止）\quad 设 $(X_0,\Delta_0)$ 为 Q-分解 dlt 对，并假设运行 MMP 得到对 $(X_i,\Delta_i)$。若 $X_{i+1}$ 由除子收缩得到，则 $\rho(X_{i+1})=\rho(X_i)-1$；若由翻转得到，则 $\rho(X_{i+1})=\rho(X_i)$。由此得到不等式
\[
\rho(X_0)\ge\rho(X_1)\ge\cdots\ge\rho(X_i)\ge\cdots
\]
其中对除子收缩为严格不等号。因此除子收缩的次数不超过 $\rho(X_0)-1$。这一论证对翻转的次数毫无说明。

引理 3.38 表明翻转至少使某一个差异系数增大，但由此给出翻转次数的一般上界似乎并不容易。然而在许多情形这是可能的。其中若干在第 6 章讨论。参见 [K+92]。\hfill$\square$
% [Tutor Note: 除子收缩使 $\rho$ 严格下降，次数受 $\rho(X_0)-1$ 控制；翻转保持 $\rho$ 不变，其终止性只能依靠差异系数的离散性，这至今仍是高维的核心困难。]

\par\medskip\noindent\textbf{3.46}\quad 尚未解决的最重要问题是 $(K+\Delta)$-翻转的存在性。已在下列情形得到证明。

\par\noindent（1）$\dim X=n$，$X$ 具有忠实 $(\C^*)^n$-作用 [Rei83a]，或具有其他某种代数群作用 [BK94]。
